\documentclass[10pt]{article}
\usepackage[a4paper,margin=16mm]{geometry}
\usepackage[no-math]{fontspec}
\setmainfont{Latin Modern Roman}
\usepackage{amsmath,amsfonts,amssymb,amsthm,mathtools,graphicx,xurl}
\usepackage[shortlabels]{enumitem}
\usepackage[unicode,hidelinks]{hyperref}
\usepackage[capitalize]{cleveref}
\allowdisplaybreaks
\setlength\emergencystretch{3em}
\setmainfont[SmallCapsFont={Latin Modern Roman Caps}]{Latin Modern Roman}
\newtheorem{thm}{Theorem}[section]
\newtheorem{lem}[thm]{Lemma}
\newtheorem{prop}[thm]{Proposition}
\newtheorem{cor}[thm]{Corollary}
\newcommand{\NN}{{\mathbb N}}
\newcommand{\er}{\mathrm{e}}
\newcommand{\cS}{\ensuremath{\mathcal{S}}}
\newcommand{\cT}{\ensuremath{\mathcal{T}}}
\newcommand{\PR}{\mathbb P}
\newcommand{\E}{{\mathbb E}\,}
\newcommand{\be}{\begin{equation}}
\newcommand{\ee}{\end{equation}}
\newcommand{\ssum}[1]{\sum_{\substack{#1}}}
\newcommand{\sprod}[1]{\prod_{\substack{#1}}}
\newcommand{\lam}{\ensuremath{\lambda}}
\renewcommand{\a}{\ensuremath{\alpha}}
\renewcommand{\b}{\ensuremath{\beta}}
\newcommand{\eps}{\ensuremath{\varepsilon}}
\renewcommand{\le}{\leqslant}
\renewcommand{\leq}{\leqslant}
\renewcommand{\ge}{\geqslant}
\renewcommand{\geq}{\geqslant}
\newcommand{\fl}[1]{{\ensuremath{\left\lfloor {#1} \right\rfloor}}}
\newcommand{\cl}[1]{{\ensuremath{\left\lceil #1 \right\rceil}}}
\renewcommand{\(}{\left(}
\renewcommand{\)}{\right)}
\newcommand{\pfrac}[2]{\left(\frac{#1}{#2}\right)}
\newcommand{\one}{\ensuremath{\mathbbm{1}}}
\newfontfamily\indicatorfont[Path=/usr/share/fonts/truetype/noto/]{NotoSansMath-Regular.ttf}
\newcommand{\finiteindicator}[1]{\begingroup\fontsize{#1}{#1}\selectfont\upshape\indicatorfont\char"1D7D9\endgroup}
\newcommand{\mathbbm}[1]{\def\finitearg{#1}\def\finitedigit{1}\ifx\finitearg\finitedigit\mathord{\mathchoice{\hbox{\finiteindicator{\the\fontdimen6\textfont0}}}{\hbox{\finiteindicator{\the\fontdimen6\textfont0}}}{\hbox{\finiteindicator{\the\fontdimen6\scriptfont0}}}{\hbox{\finiteindicator{\the\fontdimen6\scriptscriptfont0}}}}\else\PackageError{finiteglyph}{Only printed double-struck 1 is supported}{}\fi}

\begin{document}\begin{center}\Large Problem 2657: The uniform no-small-cycle probability and Brun--Hooley sieve error\end{center}
\noindent\textbf{Authors:} Kevin Ford\par
\noindent\textbf{Work/version:} Cycle Type of Random Permutations: a Toolkit. Exact arXiv 2104.12019v3; Discrete Analysis 2022:9, 36 pp., DOI 10.19086/da.38090\par
\noindent\textbf{Difficulty:} H1: After elementary inclusion--exclusion, harmonic sums, Stirling bounds and exact binomial moments, the human proof still designs a logarithmic hierarchy of cycle-length intervals and unequal even truncation orders under a shared total-length budget. Its main product and separate first-failure error product must be evaluated simultaneously and optimized through factorial tails to achieve the u log u minus u log log log u exponent uniformly. That multilevel sieve construction exceeds ordinary single-variable inclusion--exclusion, so the independent original no-small-cycle outcome is conservatively H1.\par
\noindent\textbf{Selection and ordinary prerequisites:} Only original Theorem 1.15 is selected: $1\le m\le n$, a uniformly random permutation of $S_n$, $u=n/m$, exact $e^{-H_m}$ Poisson prediction and the original uniform two-regime g(u) error. Original uniform sampling, cycle-count, harmonic-sum and indicator definitions are included. The full joint-moment proof (Theorem 1.3 and Lemma 3.1), Cauchy formula and its proof, the complete local-law Lemma 4.1/Theorem 1.5 with Corollary 1.6 for the small-u regime, and the entire Brun--Hooley hierarchy/product/error/truncation/factorial-tail proof are editable. Ordinary prerequisites are elementary multinomial counting, inclusion--exclusion (Lemma 2.3), harmonic-sum bounds (Lemma 2.1) and Stirling formula (Lemma 2.2), all precisely stated. Ford--Halberstam [34] is method attribution: the full adapted sieve construction is printed here, not deferred to that paper. Authentic source literals, including the $m_k$ factorial exponent at 1154, switched $m_{j,r}/m_{r,j}$ subscripts at 1201--1205, the $m_t/H(I_t)$ form in the local-law sum at 1307, and the unpowered $H(I_j)$ factor in the intermediate product at 1655, are retained without a mathematical verdict or formula repair. Every original selected index and equation number is restored. Joint Poisson and grouped/regime variants share this sieve core and are uncounted.\par
\noindent\textbf{Attribution and changes:} Public carrier is a strictly checked transcription of the exact printed CC BY 3.0 article. The newly acquired nonexclusive native source is only a private read-only unexecuted aid; its archive, classes, full preamble, unprinted comments and inactive body are excluded. Whole exact published original PDF is retained with its actual CC BY 3.0 first-page annotation. The historical metadata claim inferred CC4 from a generic publisher policy; the actual retained exact PDF CC3 annotation is used here. Original human mathematical tokens are checked against printed pages; presentation changes are independent installed standard typography, finite exact indicator glyph, selection and restored original numbering.\par
\medskip\hrule\medskip
\section{Introduction}
Let $\sigma$ denote a random permutation from
the symmetric group $\cS_n$, each permutation being equally
likely\footnote{Random permutations sampled
from certain other distributions have been studied, e.g. \cite{ABTbook}, but we will not discuss these here.}.
We denote by $\PR_n$ and $\E_n$ the probability and expectation
with respect to a uniform random $\sigma\in \cS_n$.
Often, the subscript $n$ will be omitted if it is clear from the context.
We denote the type (or cycle type) of $\sigma$ by
\[
(C_1(\sigma), C_2(\sigma),\ldots, C_n(\sigma)),
\]
where
$C_j(\sigma)$ is the number of cycles of length $j$ in $\sigma$.
More generally, for any subset $I$ of $[n]=\{1,\ldots,n\}$, we let $C_I(\sigma)$
be the number of cycles whose lengths lie in the set $I$.
For brevity, we write $C(\sigma)$ for the total number of
cycles in $\sigma$.
\setcounter{thm}{1}
\begin{thm}[Cauchy]\label{thm:Cauchy}
If $m_1+2m_2+\cdots+nm_n=n$, then
\[
\PR_n \big( C_1(\sigma)=m_1,\ldots,C_n(\sigma)=m_n \big) = \prod_{j=1}^n \frac{(1/j)^{m_j}}{m_j!}.
\]
\end{thm}
If  $X_1, X_2, \ldots, X_k$ are independent Poisson random variables 
with parameters $\lambda_1,\ldots,\lambda_k$, respectively, then the sum
$X_1+\cdots +X_k$ is Poisson with parameter $\lambda_1+\cdots+\lambda_k$.
Thus, for subsets $I$ of $[n]$ we should expect that $C_I(\sigma)$ will be roughly 
Poisson with parameter 
\[
H(I) := \sum_{j\in I} \frac{1}{j}.
\]
In the important special case $I=\{1,\ldots,n\}$ we set 
\[
H_n = \sum_{i=1}^n \frac{1}{i}.
\]
\subsection{Notational conventions.}
We adopt the standard Bachman-Landau, Hardy, and Vinogradov notations:
$f=O(g)$ and $f\ll g$ mean that there is a positive constant $C$
so that $|f| \le Cg$ throughout the domain of $f$.  The constant $C$ is independent of any parameters, unless specified by subscripts, e.g.
$f(x)=O_{\eps} (x^\eps)$.
Also, $f(x) \sim g(x)$ as $x\to \infty$ means $\lim_{x\to \infty} f(x)/g(x)=1$
and $f(x)=o(g(x))$ means that $\lim_{x\to\infty} f(x)/g(x)=0$.

For $\sigma\in \cS_n$, the notation
$\beta|\sigma$ means that $\beta$ is a \emph{divisor} of the permutation
$\sigma$, i.e. a product of some subset of the cycles of $\sigma$.
$|\beta|$ is the size (length) of $\beta$.

$\one(S)$ is the indicator function of statement $S$; $\displaystyle \one(S) = 
\begin{cases} 1 & S \text{ is true} \\ 0 & S \text{ is false}. \end{cases}$

\setcounter{thm}{2}
\subsection{Binomial moments.}
A great deal of our analysis ultimately relies on
estimates for joint binomial moments of the quantities
$C_I(\sigma)$.  Recall that if $X$
is Poisson with parameter $\lambda$,
then for any non-negative integer $m$,
\[
\E \binom{X}{m} = 
\sum_{k=0}^\infty \binom{k}{m} \er^{-\lambda} \frac{\lambda^k}{k!}= \frac{\lambda^m}{m!}.
\]
We establish an analog for joint binomial moments of the statistics
$C_{I_j}(\sigma)$ for disjoint $I_1,\ldots,I_k$.

\begin{thm}\label{cycles_sets_expectation}
Let $I_1,I_2, \ldots, I_k$ be disjoint, nonempty subsets of $[n]$, and let
$m_1,\ldots,m_k$ be non-negative integers.  Then
\[
\E \binom{C_{I_1}(\sigma)}{m_1} \cdots \binom{C_{I_k}(\sigma)}{m_k} \le \prod_{j=1}^k \frac{H(I_j)^{m_j}}{m_j!},
\]
with equality if and only if $\sum_{j=1}^k m_j \max(I_j) \le n$.
\end{thm}
\subsection{Local limit theorems}
\setcounter{thm}{4}
\begin{thm}\label{cycles_sets}
Let $I_1,\ldots,I_r$ be arbitrary disjoint, nonempty subsets of $[n]$ and $m_1,\ldots,m_r\ge 0$.  Then
\[
\PR\big( C_{I_1}(\sigma)=m_1,\ldots,C_{I_r}(\sigma)=m_r \big) \le \frac{\er^{H_n}}{n} \, \prod_{j=1}^r \Bigg( \frac{H(I_j)^{m_j}}{m_j!} \er^{-H(I_j)}\Bigg) \cdot \(\eps+\frac{m_1}{H(I_1)}+\cdots+\frac{m_r}{H(I_r)}\),
\]
where $\eps=0$ if $[n]=I_1\cup \cdots \cup I_r$ and $\eps=1$
otherwise.
\end{thm}
\begin{cor}\label{cycles-single-set}
For any $I\subset [n]$ and $m\ge 0$,
\[
\PR\( C_{I}(\sigma)=m \) \le \frac{\er^{H_n-H(I)}}{n} \cdot \frac{H(I)^m}{m!}
\Bigg( \one(I\ne [n]) + \frac{m}{H(I)} \Bigg).
\]
In particular,
\[
\PR( C_I(\sigma)=0 ) \le \frac{\er^{H_n-H(I)}}{n} = \er^{\gamma-H(I)}(1+O(1/n)).
\]
\end{cor}
\setcounter{subsection}{4}\setcounter{thm}{14}
\subsection{Permutations without small cycles.}
%
%

Sharp bounds on $\PR(C_{[m]}(\sigma)=0)$ are a
key to establishing the Poisson model.
The model predicts that $\PR(C_{[m]}(\sigma)=0)$
should be about $\er^{-H_m}$,
and Corollary \ref{cycles-single-set} contains an upper bound
close to this.
 This cannot be expected to hold for large $m$, for example 
 $\PR(C_{[m]}(\sigma)=0)=1/n$ if $m \ge n/2$ since a permutation lacking cycles of length at most $m$ must be a single $n$-cycle.
  In fact, when $n/m$ is small, there is an asymptotic formula
 $\PR(C_{[m]}(\sigma)=0) \sim \omega(n/m)/m$ ($n\to\infty$, $m\to\infty$) where $\omega$ is Buchstab's function
 and $\omega(u)\to \er^{-\gamma}$ as $u\to\infty$ \cite[Theorem 5]{granville}.  This is analogous to the problem
 of counting integers $n\le x$ with no prime factor $\le x^{1/u}$
 (see \cite[Ch. III.6]{Tenbook}).

Our focus  is to prove that
$\PR(C_{[m]}(\sigma)=0)$ is very close to $\er^{-H_m}$ when 
$n/m$ is large.

 \begin{thm}\label{nosmallcycles2}
 Let $1\le m\le n$.  Then
 \[
 \PR(C_{[m]}(\sigma)=0) = \er^{-H_{m}} \(1 + O(\er^{-g(n/m)}) \),
 \]
 where $g(x)=0$ for $1\le x\le 20$ and for $x > 20$,
 \[
g(x) = x\log x - x\log\log\log x + O(x).
 \]
 \end{thm}
 
 Theorem \ref{nosmallcycles2} will be proved in section \ref{sec:nosmall}.
\section{Preliminaries}

The following standard bounds are stated without proof.

\begin{lem}\label{harmonic}
The harmonic sums $H_n$ satisfy

(i) $\log n \le H_n \le 1 +\log n$;

(ii) $H_n = \log n + \gamma + O(1/n)$,
where
$\gamma=0.57721566\ldots$ is Euler's constant.
\end{lem}

\begin{lem}[Stirling's formula]\label{Stirling}
We have $n! \ge (n/\er)^n$ and the asymptotic
\[
n! = \sqrt{2\pi n} (n/\er)^n (1+O(1/n)) \qquad (n\ge 1).
\]
\end{lem} 

\begin{lem}[Inclusion-exclusion]\label{incl-excl}
Let $a$ be a non-negative integer. For $0\le m\le k$,
\begin{align*}
\one(a=m) &= \sum_{r=m}^\infty (-1)^{r-m} \binom{r}{m} \binom{a}{r}
\\ &= \sum_{r=m}^k (-1)^{r-m} \binom{r}{m} \binom{a}{r} +
(-1)^{k+1-m} \binom{a}{m} \binom{a-m-1}{k-m},
\end{align*}
where the final term is
 at most $\binom{a}{k+1}\binom{k+1}{m}$ in
absolute value.
\end{lem}

The final claim comes from the inequality $\binom{a-m-1}{k-m} \le \binom{a-m}{k-m+1}$.
\setcounter{equation}{12}
\section{Binomial moments}\label{sec:binom}
%
%
%

We begin by proving a special case of Theorem \ref{cycles_sets_expectation},
where each set $I_j$ is a singleton.  This is Theorem 7 in \cite{watterson}.

\begin{lem}\label{cycles}
Let $m_1,\ldots,m_n$ be non-negative integers with $m_1+2m_2+\cdots+nm_n\le n$.  Then
 \[
\E \prod_{j=1}^n \binom{C_j(\sigma)}{m_j} = \prod_{j=1}^n \frac{(1/j)^{m_j}}{m_j!}.
\]
If  $m_1+2m_2+\cdots+nm_n > n$, then the left side is zero.
\end{lem}
\begin{proof}
The second assertion is obvious, since the only way for the product on the left to be positive is for the sum of the cycle lengths to exceed $n$. Now assume that
$m_1+2m_2+\cdots+nm_n\le n$. 
 The number of ways of choosing from $[n]$ a disjoint collection 
  of $m_1$ $1-$element sets, $m_2$  $2-$element sets,
$\ldots$, $m_n$ $n-$element sets is equal to
\[
 \binom{n}{\underbrace{1 \cdots 1}_{m_1} \underbrace{2 \cdots 2}_{m_2} \cdots \underbrace{n\cdots n}_{m_n} \, t} 
 \frac{1}{m_1! \cdots m_n!} = \frac{n!/t!}{\prod_{j=1}^n (j!)^{m_j} m_j!},
 \]
where $t=n-(m_1+2m_2+\cdots+nm_n)$.
A $k$-element set may be arranged into a cycle in $(k-1)!$ ways.
Thus, 
the number of ways to arrange the elements of these sets into cycles is $(0!)^{m_1} (1!)^{m_2} \cdots (n-1)!^{m_k}$.
Finally, the $t$ elements not used in any of these cycles may be permuted in $t!$ ways.
\end{proof}
\begin{proof}[Proof of Theorem \ref{thm:Cauchy} (Cauchy's Theorem)] 
Apply Lemma \ref{cycles}, noting that $\binom{C_j(\sigma)}{m_j}\ne 0$
for all $j$ if and only if $C_j(\sigma)=m_j$ for every $j$.
\end{proof}
\begin{proof}[Proof of Theorem \ref{cycles_sets_expectation}]
Consider a set $A$ of size $C_{I_j}(\sigma)$,
and partition $A$ into subsets $A_r$, where
$|A_r|=C_r(\sigma)$ for $r\in I_j$.  Then
\[
\prod_{j=1}^k\binom{C_{I_j}(\sigma)}{m_j} = \sum_{\eqref{cycles-sets-2}}\;
\prod_{j=1}^k \prod_{r\in I_j} \binom{C_r(\sigma)}{m_{j,r}},
\]
where the summation is over tuples $(m_{j,r})_{1\le j\le k,r\in I_j}$
satisfying the system
\be\label{cycles-sets-2}
\sum_{r\in I_j} m_{j,r} = m_j \quad (1\le j\le k). 
\ee
Thus,
\be\label{cycles-sets-1}
\E \prod_{j=1}^k \binom{C_{I_j}(\sigma)}{m_j} = \sum_{\eqref{cycles-sets-2}} \E \prod_{j=1}^k \prod_{r\in I_j} \binom{C_r(\sigma)}{m_{j,r}}.
\ee
Using Lemma \ref{cycles}, the expectation on the right side of \eqref{cycles-sets-1} equals
\[
\prod_{j=1}^k \prod_{r\in I_j} \frac{(1/r)^{m_{r,j}}}{m_{r,j}!}
\]
provided that
\be\label{cycles-sets-3}
\sum_{j=1}^k \sum_{r\in I_j} r m_{r,j} \le n, 
\ee
and is zero otherwise.  

If $\sum_{j=1}^k m_j \max(I_j)\le n$, then \eqref{cycles-sets-3} will always be satisfied
as long as \eqref{cycles-sets-2} holds,
and therefore
\[
\E \prod_{j=1}^k \binom{C_{I_j}(\sigma)}{m_j} =
 \prod_{j=1}^k \sum_{\eqref{cycles-sets-2}} \prod_{r\in I_j} \frac{(1/r)^{m_{r,j}}}{m_{r,j}!} = \prod_{j=1}^k \frac{H(I_j)^{m_j}}{m_j!},
\]
as claimed.  On the other hand, if
 $\sum_{j=1}^k m_j \max(I_j) > n$, then there is some choice of the parameters
$(m_{j,r})$ satisfying \eqref{cycles-sets-2} but violating \eqref{cycles-sets-3}, 
and the left side is strictly less than the right side.
Specifically, we may take $m_{j,\max I_j}=m_j$ for each $j$
and $m_{j,r}=0$ otherwise. 
\end{proof}
\section{Local limit theorems}\label{sec:local}
We begin with a rather complicated identity for the joint distribution
of the quantities $C_{I_i}$.

\begin{lem}\label{lem:CIjmj}
Let $I_1,\ldots,I_r$ be disjoint subsets of $[n]$ and $m_1,\ldots,m_r$
be non-negative integers.  Denote $I_0 = [n] \setminus (I_1\cup \cdots \cup I_r)$.
Let $\cT$ be the set of indices $i$ with $m_i>0$, together
with the number 0 if $I_0$ is nonempty.  Then
\[
  \PR\big( C_{I_j}(\sigma)=m_j\; (1\le j\le r) \big)
  = \frac{1}{n}\sum_{t\in \cT} \sum_{h \in I_t} 
  \ssum{b_1,\dots, b_n \geq 0 \\ b_1 + 2b_2 + \cdots + nb_n = n - h \\ \sum_{i\in I_j} b_i =m_j-\one(t=j),\  (1\le j\le r) } 
  \prod_{i=1}^n \frac{(1/i)^{b_i}}{b_i!}.
\]
\end{lem}

\begin{proof}
 Evidently
\[
  n \# \{\sigma \in \cS_n:  C_{I_1}(\sigma)=m_1,\ldots,C_{I_r}(\sigma)=m_r\} = \ssum{\sigma \in \cS_n \\ C_{I_j}(\sigma)=m_j \; (1\le j\le r) } \;\; \ssum{\alpha|\sigma \\ \a \text{ a cycle}} |\a|.
\]
Write $\sigma=\a \beta$ and let $h=|\a|$.   
 Thus, for some $t\in \cT$,  we have $|\a|=h\in I_t$
 and
\[
(C_{I_1}(\beta),\ldots,C_{I_r}(\beta))=(m_1-\one(t=1),\ldots,m_r-\one(t=r)).
\]
It is permissible to think of $\beta\in \cS_{n-h}$ and thus
\begin{align*}
  n  \# \{\sigma \in \cS_n: C_{I_1}(\sigma)=m_1,\ldots,C_{I_r}(\sigma)=m_r\}
  &= \sum_{t\in \cT} \sum_{h\in I_t} \; \ssum{\a\in\cS_n, |\a|=h \\ \a\text{ a cycle}} h \ssum{\beta\in \cS_{n-h}\\ C_{I_i}(\beta)=m_i-\one(t=i), (1\le i\le r)} 1  \\
  &=  \sum_{t\in \cT} \sum_{h\in I_t} \frac{n!}{(n-h)!} \ssum{\beta\in \cS_{n-h} \\ C_{I_i}(\beta)=m_i-\one(t=i), (1\le i\le r) } 1.
\end{align*}
Now subdivide the sum according the cycle type $(b_1,\ldots,b_{n})$ of
the permutation $\beta$, use Cauchy's formula (Thm. \ref{thm:Cauchy}) to 
count such permutations for each type, and divide by $n$.  The desired identity follows.
\end{proof}




\begin{proof}[Proof of Theorem \ref{cycles_sets}]
The right side in Lemma \ref{lem:CIjmj} is at most
 \[
 \frac{1}{n} \sum_{t\in \cT}
  \ssum{b_1,\dots, b_n \geq 0 \\ \sum_{i\in I_j} b_i =m_j-\one(t=j) \; (1\le j\le r)  } \prod_i \frac{(1/i)^{b_i}}{b_i!} := \frac{Y}{n},
\]
  say.   By the multinomial theorem,
  \begin{align*}
 Y & =    \sum_{t\in \cT} \ssum{b_i\ge 0 \; (i\in I_1\cup \cdots \cup I_r) \\ \sum_{i\in I_j} b_i = m_j-\one(t=j) \; (1\le j\le r) } 
  \frac{1}{\prod_{i\in I_1\cup \cdots \cup I_r} b_i! i^{b_i}} \sum_{b_i\ge 0 \; (i\in I_0)} \frac1{\prod_{i\in I_0} b_i! i^{b_i}}\\
  &=  \sum_{t\in \cT} \frac{m_t}{H(I_t)}\prod_{j=1}^r \frac{H(I_j)^{m_j}}{m_j!} \er^{H(I_0)}.
\end{align*}
 The claimed bound now follows from
$H(I_0) = H_n - H(I_1)-\cdots-H(I_r)$.
\end{proof}
\setcounter{section}{5}\setcounter{equation}{20}
\section{Permutations without small cycles}\label{sec:nosmall}
%
%


\begin{proof}[Proof of Theorem \ref{nosmallcycles2}]
Our proof is based on the Brun-Hooley sieve \cite{BrunHooley}
from number theory.
Let $K \ge \er^{10}$ be fixed and sufficiently large and let $u=n/m$.  If $u \le K$, then
Corollary \ref{cycles-single-set} implies that $\PR(C_{[m]}(\sigma)=0)\ll 1/m$ and 
the conclusion follows.  Now assume that $u > K$ and let 
\[
D = \log u,
\]
so that $D\ge 10$.
Partition $[m]$ into intervals $I_j=[z_j,z_{j-1})\cap \NN$, where $z_j=m/D^j$,
$0\le j\le \cl{\frac{\log m}{\log D}}=t$. 
Let $k_1,\ldots,k_t$ be positive, even integers, subject to
\be\label{kj-cond}
k_j \ge 10\log D \;\; (1\le j\le t), \qquad \sum_{j=1}^t \frac{(k_j+1) m}{D^{j-1}} \le n.
\ee
With $\sigma$ fixed, let
\[
x_j = \one(C_{I_j}(\sigma)=0), \quad y_j = \sum_{r=0}^{k_j} (-1)^r \binom{C_{I_j}(\sigma)}{r}.
\]
By Lemma  \ref{incl-excl}, we have
\[
0\le y_j - x_j = \binom{C_{I_j}(\sigma)-1}{k_j} \le \binom{C_{I_j}(\sigma)}{k_j+1}.
\]
Using the elementary inequality
 \[
x_1\cdots x_t \ge y_1\cdots y_t - \sum_{\ell=1}^t (y_\ell-x_\ell) \sprod{j=1 \\ j\ne \ell}^t y_j,  
 \]
 together with $\PR(C_{[m]}(\sigma)=0) =  \E x_1\cdots x_t$, we thus obtain
\be\label{small-ME}
M - E \le \PR(C_{[m]}(\sigma)=0) \le M,
\ee
where
\[
M=\E y_1\cdots y_t, \qquad E = \E \sum_{\ell=1}^t \binom{C_{I_\ell}(\sigma)}{k_\ell+1} \prod_{j\ne \ell} y_j.
\]
The condition \eqref{kj-cond} implies that 
\be\label{kjmaxIj-2}
\sum_j (k_j+1) \max I_j \le n.
\ee
  Thus, by
Theorem \ref{cycles_sets_expectation},
\begin{align*}
M &= \ssum{r_1,\ldots,r_t \\ 0\le r_j\le k_j (1\le j\le t)} 
(-1)^{r_1+\cdots+r_t} \E \binom{C_{I_1}(\sigma)}{r_1} \cdots \binom{C_{I_t}(\sigma)}{r_t} \\
&=  \ssum{r_1,\ldots,r_t \\ 0\le r_j\le k_j (1\le j\le t)} 
(-1)^{r_1+\cdots+r_t} \prod_{j=1}^t \frac{H(I_j)}{r_j!} 
= \prod_{j=1}^t \Bigg( \sum_{r_j=0}^{k_j} \frac{(-H(I_j))^{r_j}}{r_j!}    \Bigg).
\end{align*}
Since $H(I_j)=\log D+O(1)$ for every $j$, 
and recalling \eqref{kj-cond}, we have 
\begin{align}
\sum_{r_j=0}^{k_j} \frac{(-H(I_j))^{r_j}}{r_j!} &= \er^{-H(I_j)} +
 O\pfrac{H(I_j)^{k_j+1}}{(k_j+1)!} \notag \\
&=\er^{-H(I_j)} \bigg( 1 + O \pfrac{D(\log D+O(1))^{k_j+1}}{(k_j+1)!}  \bigg) \notag
\\ &=\er^{-H(I_j)} \exp \Bigg[ O \pfrac{D(\log D+O(1))^{k_j+1}}{(k_j+1)!}  \Bigg].
\label{nosmall-yj}
\end{align}
Hence, the main term satisfies
\be\label{small-M}
M = \er^{-H_m} \exp\left[ O\(\sum_{j=1}^t \frac{D(\log D+O(1))^{k_j+1}}{(k_j+1)!}\) \right].
\ee
Similarly, using \eqref{kjmaxIj-2},  the error term satisfies
\begin{align*}
E &=  \sum_{\ell=1}^t \ssum{r_j (j\ne \ell) \\ 0\le r_j\le k_j (j\ne \ell))} 
(-1)^{\sum_{j\ne \ell} r_j} \E \binom{C_{I_\ell}(\sigma)}{k_\ell+1} 
\prod_{j\ne \ell} \binom{C_{I_j}(\sigma)}{r_j} \\
&=  \sum_{\ell=1}^t
 \frac{H(I_\ell)^{k_\ell+1}}{(k_\ell+1)!} \prod_{j\ne \ell}
\Bigg( \sum_{r_j=0}^{k_j} \frac{(-H(I_j))^{r_j}}{r_j!}  \Bigg) .
\end{align*}
Hence, by \eqref{nosmall-yj},
\be\label{small-E}
E=\sum_{\ell=1}^t \frac{\er^{H(I_\ell)}(\log D+O(1))^{k_\ell+1}}
{(k_\ell+1)!} \er^{-H_m} \exp\left[ O\(\sum_{j=1}^t \frac{D(\log D+O(1))^{k_j+1}}{(k_j+1)!}\) \right].
\ee
We now take
\[
k_j=k_1+2(j-1) \quad (j\ge 1), \;\; k_1=2\fl{\frac{D-1}{D} \cdot \frac{u}{2}}-6,
\]
and readily verify that the conditions \eqref{kj-cond} hold if
$K$ is large enough.
Thus, by Stirling's formula,
\begin{align*}
\sum_{j=1}^t \frac{D(\log D+O(1))^{k_j+1}}{(k_j+1)!} 
&\ll \frac{D(\log D+O(1))^{k_1+1}}{(k_1+1)!} \\
&\le \er^{-u\log u + u\log\log\log u + O(u)}
\end{align*}
and likewise
\[
\sum_{\ell=1}^t \frac{\er^{H(I_\ell)}(\log D+O(1))^{k_\ell+1}}
{(k_\ell+1)!} \le  \er^{-u\log u + u\log\log\log u + O(u)}.
\]
Inserting these last two bounds into 
\eqref{small-M} and \eqref{small-E}, and recalling \eqref{small-ME},
the proof is complete.
\end{proof}
\begin{thebibliography}{99}

\bibitem{equalorders}
{\sc H. Acan, C. Burnette, S. Eberhard, E. Schmutz and J. Thomas}.
\emph{Permutations with equal orders},
Combin. Probab. Comput. {\bf 30} (2021), no. 5, 800--810. 
%

\bibitem{Ash}  {\sc R. B. Ash}. \emph{Information theory}.
  Corrected reprint of the 1965 original.
  Dover Publications, Inc., New York, 1990. xii+339 pp.

\bibitem{ABT}
{\sc R.~Arratia, A.~D. Barbour, and S.~Tavar\'{e}}. {\em On random polynomials
  over finite fields}, Math. Proc. Cambridge Philos. Soc., 114 (1993),
  pp.~347--368.

\bibitem{ABTbook}
{\sc R.~Arratia, A.~D. Barbour, and S.~Tavar\'{e}}.
Logarithmic Combinatorial Structures: A Probabilistic Approach, EMS Monogr. Math., EMS Publishing House, Z\" urich, 2003.

\bibitem{AT92}
{\sc R. Arratia and S. Tavar\'e}.
\newblock The cycle structure of random permutations.
\newblock {\em Ann. Probab.} 20(3) (1992), 1567--1591.
%

\bibitem{AT94}
{\sc R. Arratia and S. Tavar\'e}.
\emph{Independent process approximations for random combinatorial structures.}
Adv. Math. {\bf 104} (1994), no. 1, 90--154. 
%
%

\bibitem{BGHP}
{\sc J. Bamberg, S.P. Glasby, S. Harper, and C.~E.~Praeger}.
\emph{Permutations with orders coprime to a given integer.}
Electronic J. Combinatorics {\bf 27}, 1, (2020), 14 pp.
 
 \bibitem{BSK}
{\sc L. Bary-Soroker and G. Kozma.}
\emph{Irreducible polynomials of bounded height.}
Duke Math. J. {\bf 169} (2020), no. 4, 579--598. 

\bibitem{BSKK}
{\sc L. Bary-Soroker, G. Kozma and D. Koukoulopoulos}.
{\it Irreducibility of random polynomials: general measures.}
preprint. \texttt{arXiv:2007.14567}.

 
 \bibitem{BLGNPS}
 {\sc R. Beals, C. R. Leedham-Green, A. C. Niemeyer, C.~E.
Praeger, and {\' A}. Seress. }
\emph{Permutations with restricted cycle structure and
an algorithmic application.} 
Combin. Probab. Comput. {\bf 11}, (5), (2002), 447--464.
%

\bibitem{debruijn}
{\sc N. G. de Bruijn}.
\emph{The asymptotic behaviour of a function occurring in the theory of primes}.
J. Indian Math. Soc. (N.S.) {\bf 15} (1951), 25--32. 


\bibitem{camkan}
{\sc P.~J. Cameron and W.~M. Kantor}.
\newblock {\em Random permutations: some group-theoretic aspects.}
\newblock { Combin. Probab. Comput.}, 2(3):257--262, 1993.

\bibitem{CHM}
{\sc S.~Chowla, I.~N. Herstein, and W.~K. Moore}.
\newblock On recursions connected with symmetric groups. {I}.
\newblock {\em Canad. J. Math.}, 3:328--334, 1951.
%
%
%
%
%
%
%
%

\bibitem{dfg08}
{\sc P.~Diaconis, J.~Fulman, and R.~Guralnick}.
\newblock On fixed points of permutations.
\newblock {\em J. Algebraic Combin.}, 28(1):189--218, 2008.

\bibitem{dickman}
{\sc K. Dickman}.  \emph{On the Frequency of Numbers Containing Prime Factors of a Certain Relative Magnitude}. Arkiv f\"or Mat., Astron. och Fys. 22A, 1--14, 1930.


\bibitem{dixon92} {\sc J.~Dixon}. \emph{Random sets which invariably generate the symmetric group}, Discrete Math. \textbf
{105} (1992), 25--39.


\bibitem{EFG1}
{\sc S. Eberhard, K. Ford, and B. Green}.
\newblock Permutations fixing a {$k$}-set.
\newblock {\em Int. Math. Res. Not. IMRN}, (21):6713--6731, 2016.

\bibitem{EFG2}
{\sc S. Eberhard, K. Ford, and B. Green}.
\newblock \emph{Invariable generation of the symmetric group.}
\newblock Duke Math. J. {\bf 166} (2017), no. 8, 1573-1590.

\bibitem{EFK}
{\sc S.~Eberhard, K.~Ford and D. Koukoulopoulos},
\emph{Permutations contained in transitive subgroups},
Discrete Analysis {\bf 2016: 12}, 34 pages. 

\bibitem{Erdos55}
{\sc P.~Erd{\H{o}}s}.
\newblock Some remarks on number theory.
\newblock {\em Riveon Lematematika}, 9:45--48, 1955.
\newblock (Hebrew. English summary).

\bibitem{Erdos60}
{\sc P.~Erd{\H{o}}s}.
\newblock {A}n asymptotic inequality in the theory of numbers.
\newblock {\em Vestnik Leningrad. Univ.}, 15(13):41--49, 1960.
\newblock (Russian).

\bibitem{ET1}
{\sc P. Erd{\H o}s and P. Tur{\' a}n}.
\emph{On some problems of a statistical group-theory. I.} Z. Wahrscheinlichkeitstheorie und Verw. Gebiete  {\bf 4}  (1965), 175--186.
%
%
 
 
 \bibitem{ET2}
{\sc P. Erd{\H o}s and P. Tur{\' a}n}.
 \emph{On some problems of a statistical group-theory. II},
 Acta Math. Acad. Sci. Hungar.  {\bf 18}  (1967), 151--163;
 
  \bibitem{ET3}
 {\sc P. Erd{\H o}s and P. Tur{\' a}n}.
 \emph{On some problems of a statistical group-theory. III},
 Acta Math. Acad. Sci. Hungar.  {\bf 18}  (1967), 309--320;
 
  \bibitem{ET4}
 {\sc P. Erd{\H o}s and P. Tur{\' a}n}.
 \emph{On some problems of a statistical group-theory. IV},
 Acta Math. Acad. Sci. Hungar.  {\bf 19}  (1968), 413--435.
		
 \bibitem{ET5}
 {\sc P. Erd{\H o}s and P. Tur{\' a}n}.
 \emph{On some problems of a statistical group-theory. V},
 Period. Math. Hungar.  {\bf 1}  (1971),  no. 1, 5--13.

 \bibitem{ET6}
 {\sc P. Erd{\H o}s and P. Tur{\' a}n}.
 \emph{On some problems of a statistical group-theory. VI},
 J. Indian Math. Soc.  {\bf 34}  (1971),  no. 3-4, 175--192.
		
 \bibitem{ET7}
 {\sc P. Erd{\H o}s and P. Tur{\' a}n}.
 \emph{On some problems of a statistical group-theory. VII},
 Period. Math. Hungar.  {\bf 2}  (1972), 149--163.
		
%
%
%
%
%

\bibitem{FlSe}
{\sc P. Flajolet and R. Sedgewick}.
\emph{Analytic combinatorics}. Cambridge University Press, Cambridge, 2009. xiv+810 pp. 


\bibitem{F08} {\sc K. Ford}.
\newblock The distribution of integers with a divisor in a given interval.
\newblock {\em Ann. of Math. (2)}, 168(2):367--433, 2008.

\bibitem{F08b}
 {\sc K. Ford}.
\newblock Integers with a divisor in {$(y,2y]$}.
\newblock In {\em Anatomy of integers}, volume~46 of {\em CRM Proc. Lecture
  Notes}, pages 65--80. Amer. Math. Soc., Providence, RI, 2008.


\bibitem{primepoisson}
 {\sc K. Ford}.
{\it Joint Poisson distribution of prime factors in sets},
Math. Proc. Cambridge Phil. Soc., to appear.


\bibitem{FGK}
{\sc K. Ford, B. Green and D. Koukoulopoulos}.
 \emph{Equal sums in random sets and the concentration of divisors.}  preprint, \texttt{arXiv:1908.00378}

\bibitem{BrunHooley}
{\sc K. Ford and H.~Halberstam}.
\newblock The {B}run-{H}ooley sieve.
\newblock {\em J. Number Theory}, 81(2):335--350, 2000.


\bibitem{100prisoners}
{\sc A. G{\' a}l and P. B.  Miltersen},
\emph{The cell probe complexity of succinct data structures}, Proceedings 30th International Colloquium on Automata, Languages and Programming (ICALP), (2003),  332-344.

\bibitem{Galambos76}
{\sc J. Galambos}.
\emph{The sequences of prime divisors of integers.}
Acta Arith. {\bf 31} (1976), no. 3, 213--218. 

\bibitem{GPU}
{\sc S. P. Glasby, C. E. Praeger and W.  R. Unger}.
\emph{Most permutations power to a cycle of small prime length},
 Proc. Edinb. Math. Soc. (2) {\bf 64} (2021), no. 2, 234--246. 


\bibitem{GS}
{\sc W. M. Y. Goh and E. Schmutz}.
\emph{The expected order of a random permutation.}
Bull. London Math. Soc. {\bf 23} (1991), no. 1, 34--42. 


\bibitem{Gon44}
{\sc V.~Gontcharoff}.
\newblock Du domaine de l'analyse combinatoire.
\newblock {\em Bull. Acad. Sci. URSS S\'er. Math. [Izvestia Akad. Nauk SSSR]},
  8:3--48, 1944. (Russian).
  English translation: V. Gon{\^ c}arov,
  On the field of combinatory analysis.
Amer. Math. Soc. Transl. (2) {\bf 19}, 1962, 1--46. 

\bibitem{granville}
{\sc A. Granville}.
\newblock Cycle lengths in a permutation are typically {P}oisson.
\newblock {\em Electron. J. Combin.}, 13(1):Research Paper 107, 23, 2006.

\bibitem{Gruder}
{\sc O. Gruder}.
\emph{Zur Theorie der Zerlegung von Permutationen in Zyklen.} (German).
Ark. Mat. {\bf 2} (1952), 385--414. 

\bibitem{Divisors}
{\sc R.~R. Hall and G.~Tenenbaum}.
\newblock {\em Divisors}, volume~90 of {\em Cambridge Tracts in Mathematics}.
\newblock Cambridge University Press, Cambridge, 1988.

\bibitem{HR17} {\sc G. H. Hardy and S. Ramanujan}.
\emph{The normal number of prime factors of a number $n$},
Quart. J. Math. Oxford {\bf 48}, 76--92.

\bibitem{hildebrand}
{\sc A. Hildebrand}.
\emph{On the number of positive integers $\le x$ and free of prime factors $>y$}.
J. Number Theory {\bf 22} (1986), no. 3, 289--307.

\bibitem{Jordan}
{\sc C. Jordan}.
\emph{Sur la limite de transitivit{\' e} des groupes non altern{\' e}s}, Bull. Soc. Math. France {\bf 1} (1872/73), 40--71. (French)


\bibitem{KTP}
{\sc D. E. Knuth and L. Trabb Pardo}.
\emph{Analysis of a simple factorization algorithm}.
Theoret. Comput. Sci. {\bf 3} (1976/77), no. 3, 321--348. 


\bibitem{Kch}
{\sc V. P. Kolchin and V. P. Chistyakov}.
\emph{On the cyclic structure of random permutations.} (Russian)
Mat. Zametki {\bf 18} (1975), no. 6, 929--938. 


\bibitem{Koubook}
{\sc D. Koukoulopoulos}. \emph{The distribution of prime numbers},
Amer. Math. Soc., Graduate Studies in Math. {\bf 203}, 2019.

\bibitem{lagarias}
{\sc J. Lagarias}.
\emph{Euler's constant: Euler's work and modern developments}.
Bull. Amer. Math. Soc. (N.S.) {\bf 50} (2013), no. 4, 527--628. 

\bibitem{Landau-perm}
 {\sc E. Landau}.
\emph{\" Uber die Maximalordnung der Permutationen gegebenen Grades [On the maximal order of permutations of given degree]}, 
Arch. Math. Phys. Ser. 3, vol. 5, 1903. (German)

\bibitem{Landau} {\sc E. Landau}.
\emph{Handbuch der Lehre von der Verteilung der Primzahlen},
Chelsea, 1951.  Reprint of the 1909 original. (German)



\bibitem{LS}
{\sc S. P. Lloyd and  L. A. Shepp}.
\emph{Ordered cycle lengths in a random permutation},
Trans. Amer. Math. Soc. {\bf 121} (1966), 340--357. 

\bibitem{LP}
{\sc T.~{\L}uczak and L.~Pyber}.
\newblock On random generation of the symmetric group.
\newblock {\em Combin. Probab. Comput.}, 2(4):505--512, 1993.

\bibitem{Manst-LIL}
{\sc E. Manstavi\v{c}ius}.
\newblock Iterated logarithm laws and the cycle lengths of a random
  permutation.
\newblock In {\em Mathematics and computer science. {III}}, Trends Math., pages
  39--47. Birkh\"{a}user, Basel, 2004.

\bibitem{ManPet16}
{\sc E.  Manstavi\v{c}ius and R. Petuchovas}.
\emph{Local probabilities for random permutations without long cycles.} Electron. J. Combin. {\bf 23} (2016), no. 1, Paper 1.58, 25 pp.

\bibitem{ManPet17}
{\sc E.  Manstavi\v{c}ius and R. Petuchovas}.
\emph{Local probabilities and total variation distance for random permutations.} Ramanujan J. {\bf 43} (2017), no. 3, 679--696.


\bibitem{Massias}
{\sc J.-P. Massias}, \emph{Majoration explicite de l'ordre maximum d'un {\' e}l{\' e}ment du groupe sym{\' e}trique}, Ann. Fac. Sci. Toulouse Math. (5) {\bf 6} (1984), no. 3-4, pp. 269--281. (French)


\bibitem{MezoWang}
{\sc I. Mez\H o and C. Wang}.
\emph{Some limit theorems with respect to constrained permutations and partitions.}
Monatsh. Math. {\bf 182} (2017), no. 1, 155--164. 
%

\bibitem{McK}
{\sc E. McKemmie}.
\emph{Invariable generation of finite classical groups},
 J. Algebra {\bf 585} (2021), 592--615. 


\bibitem{MoserWyman}
{\sc L. Moser and M. Wyman}.
\emph{Asymptotic development of the Stirling numbers of the first kind.}
J. London Math. Soc. {\bf 33} (1958), 133--146.

\bibitem{Nicolas}
{\sc J.-L. Nicolas}. \emph{Sur l'ordre maximum d'un {\' e}l{\' e}ment dans le groupe $S_n$ des permutations}, Acta Arithmetica {\bf 14} (1968), 315--332. (French)

\bibitem{NP}
{\sc A. C. Niemeyer and C. E. Praeger}.
\emph{On the frequency of permutations containing a long cycle.} 
J. Algebra {\bf 300} (2006), no. 1, 289--304. 

\bibitem{NP2}
{\sc A. C. Niemeyer and C. E. Praeger}.
\emph{On the proportion of permutations of order a multiple of the degree.}
J. London Math. Soc. (2) {\bf 76} (2007), no. 3, 622--632. 

\bibitem{Norton76}
{\sc K. K. Norton}.  \emph{On the number of restricted prime factors of
an integer}, Illinois J. Math. {\bf 20}, 681--705.

\bibitem{Pet16}
{\sc R. Petuchovas.}
\newblock {\em 
Asymptotic analysis of the cyclic
structure of permutations}.
\newblock PhD thesis, Vilnius University, 2016.
\newblock \texttt{arXiv:1611.02934}.

\bibitem{Pet18}
{\sc R. Petuchovas.}
\newblock Asymptotic estimates for the number of permutations without short
  cycles.
\newblock {\em Australas. J. Combin.}, 72:1--18, 2018.


\bibitem{PPR}
{\sc R. Pemantle, Y. Peres, and I. Rivin.}
\newblock Four random permutations conjugated by an adversary generate
  {$\mathcal{S}_n$} with high probability.
\newblock {\em Random Structures Algorithms}, 49(3):409--428, 2016.

\bibitem{Rudnick19}
{\sc Z. Rudnick}.
\emph{On locally repeated values of arithmetic functions over $\mathbb{F}_q[T]$.}
With an appendix by Ron Peled.
Q. J. Math. {\bf 70} (2019), no. 2, 451--472. 
%

\bibitem{Seress}
{\sc A. Seress}.
\emph{Permutation group algorithms.}
Cambridge Tracts in Mathematics, 152. Cambridge University Press, Cambridge, 2003.

\bibitem{Tenbook}
{\sc G. Tenenbaum.}
\newblock {\em Introduction to analytic and probabilistic number theory},
  volume 163 of {\em Graduate Studies in Mathematics}.
\newblock American Mathematical Society, Providence, RI, third edition, 2015.
\newblock English. Translated from the 2008 French edition by Patrick D. F. Ion.

%
%
%
%
%
%
%

\bibitem{vershik}
{\sc A. Vershik and A. Schmidt},
\emph{Limit measures that arise in the asymptotics of symmetric groups, I.}
Theoret. Veroyatn. i Prim. {\bf 22} (1977), 72--88.  (Russian)

\bibitem{watterson}
{\sc G. A. Watterson},
\emph{The sampling theory of selectively neutral alleles.}
Advances in Appl. Probability {\bf 6} (1974), 463--488. 

\bibitem{weingartner}
{\sc A. Weingartner}.
\emph{On the degrees of polynomial divisors over finite fields.} 
Math. Proc. Cambridge Philos. Soc. {\bf 161} (2016), no. 3, 469--487. 
%



\end{thebibliography}
\end{document}
